\pagenumbering{gobble}

\name{刘波}

\contactInfo{杭州}{15 年系统级研发与安全体系建设经验}{CISSP | C/C++ / Go / Rust}{kumustone@gmail.com}{18143431580}{硕士（中国海洋大学 · 通信与信息系统）}{1985 年生}

\section{职业定位}

\sloppy
\textbf{基础安全架构与风控技术负责人}

15 年深耕基础设施研发与企业安全体系建设，主张“以系统级工程能力解决复杂的业务安全挑战”。主导搭建过企业级风控中台、API 安全与 WAF 防护体系、敏感数据识别/脱敏引擎及主机审计与合规平台，形成覆盖接入层、应用层、数据层与主机层的纵深防御。擅长将安全能力抽象为可复用的产品模块与基建底座，在高并发与强对抗场景下平衡策略复杂度与业务延迟。

\vspace{0.10cm}

\section{核心能力}
\begin{itemize}[parsep=0.2ex]
    \setstretch{1.05}
  \item \textbf{企业级安全体系构建}: 主导搭建覆盖注册/登录/交易/支付核心链路的实时风险决策体系，打通“识别→处置→审计→复盘”风险闭环；建立统一安全能力接入框架，构建纵深防御架构（网关/应用/数据/主机）。
  \item \textbf{安全基建与边界防护}: 深谙 L7 代理模型与网关架构，主导 WAF 与 API 安全能力的产品化抽象；精通 Nginx 模块化开发与 Rust（Pingora）异步模型重构，设计插件化规则引擎与灰度发布机制，实现多业务统一接入与能力复用。
  \item \textbf{数据安全与合规治理}: 自研敏感数据识别与脱敏底座（覆盖 80+ 数据类型）；落地等保三级测评与关键链路审计能力，基于自研 HIDS Agent 梳理主机风险暴露面，建立高可用审计日志体系与可追溯机制。
  \item \textbf{系统级极值优化与排障}: 精通 C/C++ 与 Go 语言底层原理，熟练运用应用层零拷贝、无锁队列及协议优化压榨系统极限性能；精通复杂高并发场景下的策略执行调优，以及内存泄漏、GC 抖动等底层故障的根因定位。
\end{itemize}

\vspace{0.10cm}

\section{工作经验}
\vspace{0.10cm}

\datedsubsection{闪捷信息科技有限公司}{安全技术专家 / 产品技术负责人}{2022.04-至今}

\thirdheading{API 安全网关产品体系 — 技术负责人}
\begin{itemize}
    \setstretch{1.05}
    \item \textbf{Nginx 网关工程重构}: 基于 Nginx 进行 C Module 二次开发，重构网关核心模块（访问控制、账号识别、自定义 JSON 解析等），代码量较原有版本减少一半，容错性与可维护性大幅提升
    \item \textbf{Nexis 下一代异步网关 (Rust/Pingora)}: 为突破 Nginx 进程模型在连接复用、阻塞任务及扩展性上的瓶颈，主导以 Pingora（Rust async）重构网关底座，插件化生命周期承载复杂安全策略；QPS +138\%，RT -63\%，扩展至 AI 网关实时审计
    \item \textbf{安全能力产品化}: 将风险规则、告警抑制、脏业务回退、灰度发布等安全策略抽象为可配置能力模块，降低交付定制成本，提升产品标准化程度
\end{itemize}

\thirdheading{核心敏感数据识别引擎 (Rust)}
\begin{itemize}
    \setstretch{1.05}
    \item \textbf{功能}: 自研高性能敏感数据识别与脱敏底座，覆盖 80+ 种敏感数据类型，支持结构化 / 半结构化 / 文件数据，统一供 API 网关、脱敏服务、Web 展示及离线 SDK 多场景调用
    \item \textbf{性能与业务价值}: 核心模块性能提升 10X，整体链路吞吐提升 1 倍以上，支撑多业务场景复用与规模化落地
\end{itemize}

\thirdheading{产品交付与客户支持}
\begin{itemize}
    \setstretch{1.05}
    \item \textbf{客户方案落地}: 深入一线理解金融与政府客户需求，从风险视角与安全体系建设角度参与客户数据安全方案制定，输出产品需求，推动功能迭代与产品亮点打造
    \item \textbf{项目交付管理}: 保障项目按时按质交付，支持 POC 测试与验收；采用敏捷开发模式，每周班车发布，持续集成，快速迭代，稳定支撑多客户并行交付
\end{itemize}

\vspace{0.10cm}

\datedsubsection{玩物得志}{安全专家 / 安全团队负责人}{2019.08-2022.03}

\thirdheading{企业级风控体系 0-1 建设}
\begin{itemize}
    \setstretch{1.05}
    \item \textbf{搭建实时风险决策中台}: 设计覆盖“流量接入、特征提取、规则演算、风险处置”的完整风控闭环，利用分布式架构支撑 10w+ QPS 风险决策。
    \item \textbf{黑产对抗与专项治理}: 围绕核心链路沉淀策略体系，主导 IM 诈骗监管、反洗钱风险排查与私域导流专项，持续压低业务资损风险。
\end{itemize}

\thirdheading{风控基础工程与治理平台}
\begin{itemize}
    \setstretch{1.05}
    \item \textbf{基础工程落地}: 搭建风控网关、业务风控、内容风控的统一接入与治理框架，沉淀通用的策略执行与编排能力。
    \item \textbf{数据分析与管理台}: 建设数据分析与运营管理台，支撑特征提取、策略灰度、效果评估与回滚闭环，提高策略迭代效率与可追溯性。
    \item \textbf{端到端落地方式}: 以“特征提取 → 数据分析 → 整体方案 → 推动落地”的方式推进跨团队协作，将安全能力产品化并稳定运行。
\end{itemize}

\thirdheading{自研反垃圾系统与合规建设}
\begin{itemize}
    \setstretch{1.05}
    \item \textbf{自研底层风控引擎}: 应用底层工程能力自研文本反垃圾引擎（DFA + 正则 + 自定义规则），引入多级缓存与规则预编译替代商业 SaaS。
    \item \textbf{降本增效}: 单机吞吐量跃升 3 倍，在保障高频文本场景极值时延的同时，年度安全成本大幅下降 60\%。
    \item \textbf{合规与审计建设}: 主导并通过等保三级测评，构建 VPN/数据库/员工登录等关键链路审计能力，推动安全标准件落地以收敛风险暴露面。
\end{itemize}

\thirdheading{安全体系建设与工程化防护}
\begin{itemize}
    \setstretch{1.05}
    \item \textbf{主机安全与告警运营}: 建设主机入侵告警与日常运营机制，推动安全意识培训与事件响应流程落地。
    \item \textbf{账号与接口安全}: 推进设备指纹、接口保护、登录与权限认证等关键链路的安全加固，减少业务风险暴露面。
    \item \textbf{对抗与评估}: 参与/组织渗透测试与人机对抗评估，将演练结果回灌到策略与平台，形成持续演进机制。
\end{itemize}

\vspace{0.10cm}

\datedsubsection{蘑菇街}{资深开发工程师}{2015.05-2019.08}

\thirdheading{全站 WAF 与流量安全体系建设}
\begin{itemize}
    \setstretch{1.05}
    \item \textbf{网关与检测体系}: 构建全站 WAF 网关与检测体系，实现规则链解耦与横向扩展，通过模块化实现实时拦截与放行。
    \item \textbf{毫秒级性能保障}: 深度优化协议层与调用链路，检测开销压缩至毫秒级，爬虫高峰超时率由 10\% 降至 0.01\%。
\end{itemize}

\thirdheading{安全开发（Golang）与主机防护体系}
\begin{itemize}
    \setstretch{1.05}
    \item \textbf{自研 HTTP 层网关与 WAF}: 参与自研网关与 WAF 的核心开发与维护，接入全站流量并保障稳定性（单机 3w+/s 量级）。
    \item \textbf{HIDS 平台建设}: 开发与维护主机入侵检测系统，覆盖命令审计、文件完整性监控、漏洞扫描与安全基线扫描等能力。
    \item \textbf{主机 ACL 与防御}: 基于 ipset/iptables 构建主机 ACL 动态防火墙，强化主机层面的细粒度访问控制与防御。
\end{itemize}

\thirdheading{安全治理与漏洞管理闭环}
\begin{itemize}
    \setstretch{1.05}
    \item \textbf{安全扫描与整改推进}: 推动主机安全扫描与中间件安全扫描，跟进业务侧整改并建立复测机制。
    \item \textbf{白盒扫描与风险收敛}: 对安全代码进行白盒扫描与问题分发，推动整改高风险缺陷并沉淀治理流程。
    \item \textbf{漏洞验证与演练}: 参与漏洞验证与内部渗透测试，将验证结果回灌到规则与基线，持续提升整体安全水位。
\end{itemize}

\thirdheading{IM 安全升级重构}
\begin{itemize}
    \setstretch{1.05}
    \item \textbf{架构重构保障}: 重构登录鉴权与会话管理体系，增强异常行为识别与滥用控制能力，支撑亿级消息规模的安全稳定运行。
    \item \textbf{基础设施沉淀}: 开发通用长连接网络库，兼顾高性能与稳定性，显著降低双端开发成本与接入风险。
\end{itemize}

\vspace{0.10cm}

\datedsubsection{早期经历：北京中交兴路 \& 华为技术 \& 千方集团}{开发工程师}{2010.07-2015.05}

\begin{itemize}
    \setstretch{1.05}
    \item \textbf{核心技术沉淀}: 主导千万级车联网接入平台架构设计，自研基于 Epoll 的半同步/半异步多线程网络库与分布式缓存；参与华为高端路由器 VRPV8 虚拟文件管理系统研发；积累了极深厚的 Linux 内核机制、C/C++ 与高并发 I/O 调优底蕴。
\end{itemize}

\vspace{0.10cm}

\section{技术栈 \& 证书}
\begin{itemize}[parsep=0.2ex]
    \setstretch{1.05}
  \item \textbf{专业资质}: \textbf{CISSP}（国际注册信息系统安全专家）, CET-6, 国家发明专利持有人, 卫星导航定位科学技术个人一等奖
  \item \textbf{核心语言}: \textbf{C/C++} (精通，10 年+), \textbf{Go} (精通), \textbf{Rust} (持续进修中，具备基本开发能力), Python
  \item \textbf{系统底座}: Nginx (C Module), Pingora, Linux Kernel (Netfilter), Docker / K8s
  \item \textbf{数据组件}: Redis, Kafka, ClickHouse, MySQL, Zookeeper
\end{itemize}

\vspace{0.10cm}

\section{教育背景}
\begin{tabularx}{\textwidth}{@{\extracolsep{\fill}}X X X X@{}}
  \textbf{中国海洋大学} & \textit{通信与信息系统} & \textit{硕士} & \textit{2007.9 - 2010.6} \\
  \textbf{怀化学院} & \textit{电子信息科学与技术} & \textit{本科} & \textit{2003.9 - 2007.6} \\
\end{tabularx}

